
\documentclass[11pt]{article}

\usepackage[margin=1in]{geometry}
\usepackage{amsmath,amssymb}
\usepackage{booktabs}
\usepackage{graphicx}
\usepackage{subcaption}
\usepackage{xcolor}
\usepackage{hyperref}

\title{Step1a: In-Sample Diagnostic of the Perturbed Fenchel--Young Surrogate}
\author{Weikang}
\date{May 2026}

% Change this path if the figures are generated elsewhere.
% Recommended formal output:
%   results/step1_runs/step1a/epsilon=<eps>/
\newcommand{\StepOneARoot}{results/step1_runs/step1a}

\begin{document}
\maketitle

\section{Purpose of Step1a}

Step1a is an in-sample diagnostic experiment designed to visualize how two
training objectives move a two-dimensional reward model through the decision
landscape of a synthetic kidney exchange problem. It is not intended to prove
held-out generalization. Its purpose is to answer the following diagnostic
question:

\begin{quote}
On a fixed batch of synthetic KEP instances, do MSE reward fitting and
decision-focused Fenchel--Young (FY) training move through different regions of
the two-dimensional synthetic-label decision-gap landscape?
\end{quote}

The experiment is useful because it makes the FY surrogate visually inspectable:
we can plot both training trajectories in the same parameter plane and evaluate
each point by the downstream KEP objective gap under the synthetic label weights.

\paragraph{Important language.}
Throughout this report, we use \emph{synthetic-label decision gap} or
\emph{oracle objective gap}. We avoid the phrase ``true regret'', because the
evaluation weights are synthetic labels rather than real clinical outcomes.

\section{Dataset and Prediction Model}

The current Step1a diagnostic uses the controlled noisy-linear synthetic reward
dataset:
\[
    \texttt{dataset/processed/step1\_noisy\_linear\_sigma010}.
\]
Each edge \(e\) has a synthetic label weight
\[
    w^{\mathrm{syn}}_e
    =
    \max\left\{0,\,
    (10u_e + 5c_e)(1+\xi_e)
    \right\},
\]
where \(u_e\) is the edge utility, \(c_e\) is the recipient cPRA feature, and the
multiplicative noise scale is \(\sigma=0.10\).

The learned reward model is the two-feature linear probe
\[
    \hat w_e(\theta)
    =
    \theta_1 u_e + \theta_2 c_e,
\]
with
\[
    \theta_1 \leftrightarrow \text{utility}, \qquad
    \theta_2 \leftrightarrow \text{recipient cPRA}.
\]
The vector \([10,5]\) is the \emph{noise-free linear coefficient}; it should not
be called the true parameter, because the labels include multiplicative
edge-level noise.

\section{Training Objectives}

\subsection{MSE reward-fitting baseline}

The two-stage baseline trains \(\theta\) by reward fitting:
\[
    \min_{\theta}
    \frac{1}{N}\sum_{i=1}^N
    \frac{1}{|E_i|}
    \sum_{e\in E_i}
    \left(
        \hat w^i_e(\theta) - w^{i,\mathrm{syn}}_e
    \right)^2.
\]
This objective tries to recover the synthetic reward labels. In this controlled
setting, it is expected to move toward the clean signal direction \([10,5]\).

\subsection{Decision-focused FY training}

FY training updates the same two-dimensional model using the perturbed
optimizer surrogate. For a graph instance, let
\[
    y^\star = S(w^{\mathrm{syn}})
\]
be the KEP solution optimized under the synthetic label weights. The FY
objective used in this diagnostic is, up to constants independent of \(\theta\),
\[
    L_{\mathrm{FY}}(\theta)
    \approx
    \mathbb{E}_{z}
    \left[
        \max_{y\in\mathcal{Y}}
        \langle \hat w(\theta)+z, y\rangle
    \right]
    -
    \langle \hat w(\theta), y^\star\rangle.
\]
The corresponding gradient estimate has the form
\[
    \nabla_{\theta} L_{\mathrm{FY}}(\theta)
    \approx
    X^\top
    \left(
        \bar y_\epsilon(\hat w(\theta)) - y^\star
    \right),
\]
where \(\bar y_\epsilon\) is the average KEP solution under perturbed predicted
weights.

\section{Evaluation Metric}

For each parameter vector \(\theta\), Step1a evaluates the downstream decision
quality by the synthetic-label decision gap:
\[
    G(\theta)
    =
    \frac{1}{N}
    \sum_{i=1}^N
    \left[
        (w^{i,\mathrm{syn}})^\top y_i^\star
        -
        (w^{i,\mathrm{syn}})^\top
        S(\hat w^i(\theta))
    \right].
\]
Here \(y_i^\star = S(w^{i,\mathrm{syn}})\) is the oracle KEP solution under the
synthetic label weights, and \(S(\hat w^i(\theta))\) is the KEP solution induced
by the predicted weights. Smaller \(G(\theta)\) is better.

\section{Experimental Configuration}

The historical completed Step1a diagnostic used the following configuration.

\begin{table}[h]
\centering
\caption{Step1a diagnostic configuration.}
\label{tab:step1a-config}
\begin{tabular}{ll}
\toprule
Item & Value \\
\midrule
Dataset & \texttt{dataset/processed/step1\_noisy\_linear\_sigma010} \\
Number of graph instances & 100 \\
Training epochs & 1000 \\
Seed & 42 \\
Initialization & \((1.6236,\,3.3521)\) \\
MSE learning rate & 0.1 \\
FY learning rate & 0.1 \\
FY perturbation samples \(M\) & 16 \\
Grid size for landscape & \(25\times25\) \\
Epsilon values & \(0.2,\ 1.0\) \\
\bottomrule
\end{tabular}
\end{table}

\section{Results: Summary Table}

\begin{table}[h]
\centering
\caption{Historical Step1a results on the noisy-linear diagnostic. The decision
gap is the average oracle objective gap under synthetic labels.}
\label{tab:step1a-results}
\begin{tabular}{lccccc}
\toprule
Method & \(\epsilon\) & Final \(\theta\) & Final FY objective
& Final decision gap & Best decision gap \\
\midrule
MSE & -- & \((10.0017,\,5.0147)\) & -- & 0.8759 & 0.8759 \\
FY & 0.2 & \((2.8135,\,1.3589)\) & 0.3609 & 0.8714 & 0.7183 at epoch 42 \\
FY & 1.0 & \((13.9796,\,6.7825)\) & 1.8042 & 0.8714 & 0.8668 at epoch 781 \\
\bottomrule
\end{tabular}
\end{table}

The table suggests three diagnostic observations. First, the MSE baseline
recovers the noise-free linear coefficient direction, but does not reach zero
decision gap because the two-feature probe cannot represent the deterministic
multiplicative edge noise. Second, FY with \(\epsilon=0.2\) visits a
lower-gap region than the final MSE endpoint, but the best-gap point occurs
early. Third, larger perturbation \(\epsilon=1.0\) changes the scale and timing
of the FY trajectory.

\section{Figure-Based Analysis}

\subsection{\texorpdfstring{\(\epsilon=0.2\)}{epsilon=0.2}: trajectory on the decision-gap landscape}

\begin{figure}[h]
\centering
\includegraphics[width=\textwidth]{\StepOneARoot/epsilon=0.2/trajectory_contour.png}
\caption{\(\epsilon=0.2\). MSE reward fitting and FY decision-focused training
trajectories on the synthetic-label decision-gap landscape. The colorbar is the
average oracle objective gap. The diamond marks the noise-free linear coefficient
\([10,5]\), the star marks the final iterate, and the gold marker marks the
best-gap iterate along the trajectory.}
\label{fig:step1a-eps02-contour}
\end{figure}

\paragraph{How to explain this figure.}
The left panel shows that MSE reward fitting moves toward the noise-free linear
coefficient \([10,5]\), which is expected because the reward labels are generated
from a noisy version of that clean linear signal. The right panel shows that FY
does not need to recover \([10,5]\). Instead, it follows the geometry induced by
the KEP solver and the perturbed FY objective. For \(\epsilon=0.2\), the
trajectory visits a low decision-gap region early, which motivates checkpoint
selection rather than only using the final iterate.

\subsection{\texorpdfstring{\(\epsilon=0.2\)}{epsilon=0.2}: 3D trajectory metrics}

\begin{figure}[h]
\centering
\includegraphics[width=\textwidth]{\StepOneARoot/epsilon=0.2/trajectory_3d_metrics.png}
\caption{\(\epsilon=0.2\). Three trajectory views: MSE decision gap, FY decision
gap, and FY perturbed objective.}
\label{fig:step1a-eps02-3d}
\end{figure}

\paragraph{How to explain this figure.}
This figure separates the decision metric from the training surrogate. The FY
objective can continue decreasing even when the deterministic KEP decision gap
has already reached a good region. This is expected because the FY objective
optimizes perturbation robustness, not only the unperturbed deterministic KEP
solution.

\subsection{\texorpdfstring{\(\epsilon=0.2\)}{epsilon=0.2}: FY objective and decision gap over epochs}

\begin{figure}[h]
\centering
\includegraphics[width=\textwidth]{\StepOneARoot/epsilon=0.2/trajectory_epoch_metrics.png}
\caption{\(\epsilon=0.2\). FY perturbed objective and synthetic-label decision
gap over epochs.}
\label{fig:step1a-eps02-epoch}
\end{figure}

\paragraph{How to explain this figure.}
The best decision-gap epoch can occur before the final epoch. Therefore, Step1a
motivates validation-based model selection in Step1b: the final FY iterate is
not necessarily the best deterministic decision-quality iterate.

\subsection{\texorpdfstring{\(\epsilon=1.0\)}{epsilon=1.0}: trajectory on the decision-gap landscape}

\begin{figure}[h]
\centering
\includegraphics[width=\textwidth]{\StepOneARoot/epsilon=1.0/trajectory_contour.png}
\caption{\(\epsilon=1.0\). MSE reward fitting and FY decision-focused training
trajectories on the synthetic-label decision-gap landscape.}
\label{fig:step1a-eps10-contour}
\end{figure}

\paragraph{How to explain this figure.}
With a larger perturbation scale, FY tends to move to a larger parameter scale.
This is consistent with the fact that the model must produce decisions that are
stable under stronger perturbations. The parameter endpoint is therefore not
expected to coincide with the clean coefficient \([10,5]\).

\subsection{\texorpdfstring{\(\epsilon=1.0\)}{epsilon=1.0}: 3D trajectory metrics}

\begin{figure}[h]
\centering
\includegraphics[width=\textwidth]{\StepOneARoot/epsilon=1.0/trajectory_3d_metrics.png}
\caption{\(\epsilon=1.0\). Three trajectory views: MSE decision gap, FY decision
gap, and FY perturbed objective.}
\label{fig:step1a-eps10-3d}
\end{figure}

\paragraph{How to explain this figure.}
The larger perturbation scale changes the shape and scale of FY optimization.
The 3D view helps distinguish movement in parameter space from movement in the
decision metric and the FY surrogate objective.

\subsection{\texorpdfstring{\(\epsilon=1.0\)}{epsilon=1.0}: FY objective and decision gap over epochs}

\begin{figure}[h]
\centering
\includegraphics[width=\textwidth]{\StepOneARoot/epsilon=1.0/trajectory_epoch_metrics.png}
\caption{\(\epsilon=1.0\). FY perturbed objective and synthetic-label decision
gap over epochs.}
\label{fig:step1a-eps10-epoch}
\end{figure}

\paragraph{How to explain this figure.}
The best decision-gap epoch occurs much later than in the \(\epsilon=0.2\)
run. This shows that \(\epsilon\) affects both the scale of the trajectory and
the timing of useful decision-quality improvements.

\section{Interpretation}

Step1a supports the following qualitative conclusions.

\begin{enumerate}
    \item \textbf{MSE reward fitting follows the reward-label geometry.}
    It moves toward the noise-free linear coefficient \([10,5]\), because its
    objective is to fit edge weights.

    \item \textbf{FY training follows a decision-induced geometry.}
    It updates parameters according to differences between perturbed optimizer
    outputs and the synthetic-label oracle KEP solution.

    \item \textbf{The best deterministic decision point need not be the final FY
    endpoint.}
    For \(\epsilon=0.2\), the best observed decision gap appears early, while the
    FY objective continues to optimize its own perturbation-smoothed criterion.

    \item \textbf{The perturbation scale changes the path and parameter scale.}
    Larger \(\epsilon\) encourages larger parameter scale and different
    trajectory timing.

    \item \textbf{Step1a motivates Step1b.}
    Because the best decision-gap iterate may occur before the final iterate,
    and because this experiment is in-sample, we need Step1b to use
    train/validation/test splitting and validation-based checkpoint selection.
\end{enumerate}

\section{What Step1a Does Not Claim}

Step1a should not be interpreted as evidence of held-out generalization. It uses
the same sampled graph instances for trajectory training, metric evaluation, and
decision-gap landscape construction. Therefore, the correct claim is:

\begin{quote}
The FY surrogate can move the two-feature probe through low synthetic-label
decision-gap regions on a fixed controlled batch.
\end{quote}

The stronger claim,

\begin{quote}
FY generalizes better than MSE on unseen KEP instances,
\end{quote}

belongs to Step1b.

\section{Transition to Step1b}

Step1a reveals that the final FY iterate is not necessarily the best point in
terms of deterministic decision quality. This naturally motivates Step1b, where
we train on different numbers of graph instances, select checkpoints using
validation metrics, and evaluate the selected models on held-out test instances.

\end{document}
